\renewcommand{\Anno}{2002}

\renewcommand{\NumeroProblema}{3}

\renewcommand{\Testo}{Su una circonferenza vengono scritti quattro
numeri 1 e cinque numeri $-1$, in un ordine qualsiasi.  Al passo successivo i nove numeri vengono cambiati simultaneamente in altri nove sostituendo ad ogni numero il suo prodotto con il numero immediatamente adiacente in senso orario.  Si continua a modificare i nove numeri per i passi successivi con la stessa regola.
	
Determinare se, scegliendo opportunamente la configurazione iniziale di quattro 1 e cinque $-1$, è possibile raggiungere una
configurazione in cui i nove numeri sono tutti uguali.}

\renewcommand{\Risposta}{È impossibile.}

\renewcommand{\Aiutino}{Se ad un certo passo ci sono tutti $+1$, cosa
c'era al passo precedente?  Se ad un certo passo ci sono tutti $-1$,
cosa c'era al passo precedente?}

\renewcommand{\SoluzioneUfficiale}{\TitoloSoluzione Supponiamo che ad
una certa mossa ci siano tutti $-1$.  Allora necessariamente alla mossa precedente c'erano tutti $+1$ e $-1$ alternati, ma questo non è possibile essendo 9 un numero dispari.  Per convincercene, supponiamo di numerare da 1 a 9 (procedendo in senso orario) le varie posizioni sulla circonferenza in cui sono scritti i numeri: se nella prima posizione c'è un $+1$, allora nella seconda c'è un $-1$, nella terza c'è un $+1$, e cos\`{\i} via fino alla nona in cui viene ad esserci un $+1$, che si ritrova adiacente al $+1$ della prima posizione.  Discorso del tutto analogo vale se nella prima posizione c'è $-1$.

Questo mostra che non è possibile arrivare ad avere tutti $-1$.

Supponiamo allora che ad una certa mossa si arrivi ad avere tutti $+1$, e consideriamo la prima mossa in cui ciò accade. Alla mossa precedente o c'erano già tutti $+1$ (ma questo contraddice il fatto che stiamo considerando la prima mossa con tutti $+1$), o c'erano tutti $-1$, cosa che abbiamo appena dimostrato essere impossibile.

Se ne deduce che non è possibile raggiungere posizioni con tutti i numeri uguali.
}

\renewcommand{\SoluzioniAlternative}{
}

\renewcommand{\Osservazioni}{\TitoloOsservazione Nella soluzione
abbiamo in realtà dimostrato che l'unico modo di raggiungere prima o poi una configurazione con tutti i numeri uguali è di partire con tutti $+1$ (e fare quante mosse ci pare, che tanto non cambia mai nulla!), oppure di partire con tutti $-1$ e non fare nessuna mossa (altrimenti non si ritorna più in quella configurazione). L'informazione data nel testo (cioè che si partiva con quattro $+1$ e cinque $-1$) è quindi quasi del tutto irrilevante, in quanto serve solo per escludere che si parta con già tutti i numeri uguali.

Abbiamo anche dimostrato che una situazione analoga si ha sostituendo
il 9 con qualunque altro intero positivo dispari.}

\renewcommand{\Commenti}{
}